\section{Problema Bin Packing}

\subsection{Introducción}

El problema de empaquetado de objetos o Problema Bin Packing se produce de forma frecuente en entornos de almacenes,
transporte y mensajería. Los objetos pueden ser paquetes, cajas o pallets\footnote{Se le llama pallet o palé a la plataforma
formada con tablas de madera que se utiliza para el almacenamiento y el traslado de mercancias} que deben ser cargados en 
contenedores, furgonetas, camiones, vagones o cualquier otro medio de transporte. El objetivo es minimizar el coste asociado
al uso de los contenedores, aunque suele reducirse al mínimo de contenedores necesarios para efectuar una carga. Dependiendo
de las características de los objetos se clasifican entre

\begin{itemize}
    \item Problemas de Bin Packing de dimensión uno, para objetos de alta densidad y solo se considera el peso.
    \item Problemas de Bin Packing de dimensión dos, para carga de objetos de la misma altura.
    \item Problemas de Bin Packing de dimensión tres, para carga de objetos de baja densidad, en cuyo caso se considera el
    volumen.
\end{itemize}

\subsection{Problema unidimensional}

Sea $n$ el número de objetos que hay que empaquetar, y para cada $j=1,\dots,n$ $a_j$ es peso del objeto $j$. Se dispone de $m$
contenedores o ``bins'', y para cada $i=1,\dots,m$ $q_i$ es la capacidad o peso total que admite el contenedor $i$ y $c_i$ es
el costo por uso del contenedor $i$. Se supone que todo objeto cabe en el contenedor. Se formula el problema de la siguiente
manera

\begin{alignat}{3}
    \text{minimizar:}   &   \sum_{i=1}^m c_{i}y_{i}               \\
    \text{sujeto a:}    &   \sum_{i=1}^m x_{ij}=1 \quad j=1,\dots,n             \\
                        &   \sum_{j=1}^n a_{j}x_{ij}\leq q_iy_i \quad i=1,\dots,m \\
                        &   x_{ij} \in {0, 1} \quad i=1,\dots,m \quad j=1,\dots,n \\
                        &   y_{i} \in {0, 1} \quad i=1,\dots,m
\end{alignat}

\noindent donde las variables binarias $y_i$ indican el uso del contenedor $i$ y las $x_{ij}$ indican la asignación de objetos
a los contenedores.

Cuando los contenedores son identicos, se toma la capacidad $q_i=q$ para todos los contenedores, y la función objetivo se 
reduce a minimizar el número de contenedores usados
\begin{alignat}{3}
    \text{minimizar:}   &   \sum_{i=1}^m y_{i}
\end{alignat}

\newpage

En algunas ocasiones, el número de contenedores que pueden usarse es ilimitado. Entonces se toma $m$ lo suficietemente grande
para que el modelo tenga solución. Por ejemplo, en el caso de contenedores idénticos basta con que $m=n$, aunque siempre es
preferible un valor menor para que no haya demasidas variables. Para encontrar un valor menor vale con encontrar una solución
factible rápida y usar el número de contenedores utilizados en la solución como valor de $m$.

\subsection{Métodos de solución}

Existen numerosos algoritmos heurísticos para el Problema de Bin Packing, que permiten encontar soluciones factibles cercanas
al óptimo. Por ejemplo, puede aplicarse el algoritmo conocido como FFD (First Fit Decreasing)

\begin{itemize}
    \item \textbf{Etapa 0}. Se ordenan los objetos de forma decreciente respecto del peso. Sea $S$ la lista de objetos, $V$
    la lista de ``bins'' y $T$ la lista de los ``bins'' usados. Inicialmente $T=\emptyset$.
    \item \textbf{Etapa 1}. Elegimos el primer objeto $j$ de la lista $S$ y lo insertamos en el primer ``bin'' $i\in T$ que
    tenga capacidad residual mayor o igual que $a_j$. Si no hay ningún ``bin'' elegimos el primero $k$ de la lista $V$, lo
    añadimos al final de la lista $T$ e insertamos $j$ en $k$.
    \item \textbf{Etapa 2}. Si $S=\emptyset$ se termina, pues todos los objetos han sido ya cargados. Entonces $T$ es el
    conjunto de ``bins'' usados mientras que $V$ son los no usados. Si $S\neq\emptyset$ volvemos a la Etapa 1.
\end{itemize}

\noindent Es posible y sencillo dar una cota de la solución óptima $z$ tal que
\[
    z_1\leq z\leq z_2
\]
\noindent donde $z_2$ es la solución proporcionada por el algoritmo anterior y $z_1$ definida como
\[
    z_1= \ceil*{\frac{\sum_{j=1}^n a_j}{q}}
\]
\noindent En ocasiones resulta ser que $z_1=z_2$, en esta situación hemos encontrado el valor mínimo sin necesidad de
resolver el modelo.

\begin{ejemplo}
    Se tienen que transportar 11 objetos desde un almacén a otro. Los pesos de los objetos vienen dados en la siguiente tabla
    
    \begin{center}
        \begin{tabular}{c|ccccccccccc}
            \hline
            objeto & 1 & 2 & 3 & 4 & 5 & 6 & 7 & 8 & 9 & 10 & 11 \\ \hline
            peso & 67.3 & 504 & 121 & 130 & 375 & 234 & 73.7 & 175 & 8.5 & 225 & 50.5 \\ \hline
        \end{tabular}
    \end{center}

    Los objetos deben ser transportados en furgonetas que admiten hasta 670 Kg (es decir, $q=670$). El algoritmo FFD encuentra
    una asignación de los 11 objetos a 3 furgonetas, por tanto $z_2=3$. Como $z_1=\ceil*{\frac{\sum_{j=1}^n a_j}{q}}=
    \ceil*{\frac{1964}{670}}=\ceil*{2.93}=3$, en este caso podemos asegurar que el mínimo número necesario de furgonetas es 3.
\end{ejemplo}
